\documentclass[12pt]{article}
\usepackage[utf8]{inputenc}
\usepackage[T1]{fontenc}
\usepackage{hyperref}
\title{Wik-Epa Language Resources for K-12 Curriculum Development}
\author{Compiled by Claude Sonnet 5.5 (claude-sonnet-5-5) for Andr\'as Kornai}
\date{}
\begin{document}
\maketitle

\section{Introduction}
Wik-Epa (ISO 639-3: wie; Glottocode: wike1239) is an extinct Paman (Pama-Nyungan) language of western Cape York Peninsula, Queensland. AIATSIS records the name Wik-Ep (Y52) and notes that speakers called it Wik-Iit. Search results suggest the Wikepa were close to extinction by the post-war period, and only a few people were known to be living at the Aurukun mission in 1958. I found no current speaker count and no teaching materials.

\section{Primary Educational Resources}
No primer, textbook, school programme or classroom material for Wik-Epa turned up in my searches. A heritage curriculum would have to be built from archival and linguistic sources and from consultation with the Wik communities at Aurukun, and it would probably use English as the medium.

\section{Linguistic and Reference Materials}
\begin{itemize}
\item AIATSIS Austlang record for Wik Epa (Y52), with alternative names and spellings: \url{https://collection.aiatsis.gov.au/austlang/language/y52}
\item Wikipedia article on Wik Epa, covering the people and their territory: \url{https://en.wikipedia.org/wiki/Wik_Epa}
\item Wikipedia article on the Wik peoples, for regional context: \url{https://en.wikipedia.org/wiki/Wik_peoples}
\end{itemize}
Peter Sutton's 1978 PhD thesis on Wik society, territory and language at Cape Keerweer is the main scholarly source cited for these languages. I saw it only as a citation, with no URL, so it is not listed in the JSON.

\section{Religious and Community Resources}
No scripture translations or community language organisation specific to Wik-Epa turned up in the searches.

\section{Digital and Audio Resources}
I found no audio, video or online corpus specific to Wik-Epa.

\section{Curriculum Development Context}
The language is extinct and the sources are thin. Any curriculum would be a cultural and historical unit rather than a language-teaching programme. It could draw on the AIATSIS record and the neighbouring documented Wik languages, with guidance from Wik elders and descendants.
\end{document}